\documentclass[12pt]{article}

\usepackage{bbm}
\usepackage{fancyhdr}
\usepackage{titling}
\usepackage{enumitem}
\usepackage{amsmath}
\usepackage{dsfont}
\usepackage{amsfonts}
\usepackage{amssymb}
\usepackage{amsthm}
\usepackage{geometry}
\usepackage{tikz}
\usepackage{algorithm}
\usepackage{algpseudocode}
\geometry{verbose,tmargin=1in,bmargin=1in,lmargin=1in,rmargin=1in}
\renewcommand{\vec}[1]{\mathbf{#1}}
\newcommand{\dec}{\text{dec}}
\newcommand{\divi}{\mathbin{\textnormal{div}}}
\title{Homework 1}
\author{Satvik Prasad}
\date{}
\newtheorem{lemma}{Lemma}[section]
\fancypagestyle{plain}{%
  \fancyhf{}
  \fancyhead[L]{\theauthor}
  \fancyhead[R]{\thedate}
}
\pagestyle{plain}
\makeatletter
\def\@maketitle{%
  \newpage
  \null
  \vskip 1em%
  \begin{center}%
    \let \footnote \thanks
    {\LARGE \@title \par}%
    \vskip 1em%
  \end{center}%
  \par
\vskip 1em}
\makeatother
\usepackage{cmbright}

\newcommand{\mat}[1]{\mathbf{#1}}

% ---------- Probability theory ----------
% Core operators (optional argument-free; use as \E{X}, \Var{X}, \Cov{X}{Y})
\DeclareMathOperator{\EOp}{\mathbb{E}}
\DeclareMathOperator{\VarOp}{Var}
\DeclareMathOperator{\CovOp}{Cov}
\DeclareMathOperator{\CorrOp}{Corr}

\renewcommand{\P}[1]{\mathbb{P}\left(#1\right)}        % \P{A}
\newcommand{\E}[1]{\EOp\left[#1\right]}                % \E{X}
\newcommand{\Var}[1]{\VarOp\left(#1\right)}            % \Var{X}
\newcommand{\Cov}[2]{\CovOp\left(#1,\,#2\right)}       % \Cov{X}{Y}
\newcommand{\Corr}[2]{\CorrOp\left(#1,\,#2\right)}     % \Corr{X}{Y}
\newcommand{\SD}[1]{\operatorname{SD}\left(#1\right)}  % \SD{X}

% Subscripted versions, e.g. \Esub{\theta}{X}, \Psub{\theta}{A}
\newcommand{\Esub}[2]{\EOp_{#1}\left[#2\right]}
\newcommand{\Psub}[2]{\mathbb{P}_{#1}\left(#2\right)}

% Conditioning, indicators, independence
\newcommand{\given}{\,\middle|\,}                      % \E{X \given Y}
\newcommand{\ind}[1]{\mathbbm{1}\left\{#1\right\}}     % \ind{X > 0}
\newcommand{\indep}{\perp\!\!\!\perp}                  % X \indep Y
\newcommand{\iid}{\overset{\text{iid}}{\sim}}          % X_i \iid \Normal(0,1)

% Convergence
\newcommand{\convp}{\xrightarrow{\;p\;}}               % in probability
\newcommand{\convd}{\xrightarrow{\;d\;}}               % in distribution
\newcommand{\convas}{\xrightarrow{\text{a.s.}}}        % almost surely
\newcommand{\convLp}[1]{\xrightarrow{L^{#1}}}          % in L^p

% Probability spaces
\newcommand{\F}{\mathcal{F}}                           % sigma-algebra
\newcommand{\B}{\mathcal{B}}                           % Borel sigma-algebra
\newcommand{\pspace}{(\Omega, \F, \mathbb{P})}

% Common distributions
\newcommand{\Normal}{\mathcal{N}}
\newcommand{\Bern}{\operatorname{Bernoulli}}
\newcommand{\Bin}{\operatorname{Binomial}}
\newcommand{\Pois}{\operatorname{Poisson}}
\newcommand{\Unif}{\operatorname{Uniform}}
\newcommand{\Exp}{\operatorname{Exp}}
\newcommand{\Geom}{\operatorname{Geometric}}
\newcommand{\Gam}{\operatorname{Gamma}}
\newcommand{\Beta}{\operatorname{Beta}}

% Number sets
\newcommand{\R}{\mathbb{R}}
\newcommand{\N}{\mathbb{N}}
\newcommand{\Z}{\mathbb{Z}}


\begin{document}
\subsection{Green Circle}
\begin{enumerate}
    \item[1] Let $m = \max(a_1, \ldots, a_k).$ Then, for any integer $n \geq 1,$ we have
    $$
    (a_1^n + a_2^n + \ldots +a_k^n)^{1/n} = m((a_1/m)^n + \ldots + (a_k/m)^n)^{1/n}.
    $$
    Let us assume that we have $j \in \mathbb Z$ elements satisfying $a_i = m,$ and the remaining with $a_i < m$ since $m$ is the maximum. Then, as $n \to \infty,$ $(a_i / m)^n \to 0$ if $a_i < m,$ and $a_i /m \to 1$ otherwise.

    Thus, we have
    $$
    \lim_{n \to \infty}(a_1^n + \ldots + a_k^n)^{1/n} = m \lim_{n \to \infty} (j)^{1/n} =  m \lim_{n \to \infty} \sqrt[n]j{} = m.
    $$
    since $j \geq 1.$

    Thus, this limit $L = \max(a_1, \ldots, a_n).$

    \item[2] Consider all continuous functions $f: \mathbb R \to \mathbb R$ with $f(2x) - f(x) = x, $ and $f(0) = 1.$ I.e.,
    $$
    f(x) = \dfrac{1}{2}x + f(x/2).
    $$

    Let $a_k = f(x/2^k)$ be terms of the sequence $\{a_k\}_{k \geq 0}.$ Then, we get
    $$
    a_n = x/2^{n + 1} + a_{n + 1}.
    $$
    Then, the limit of this series is 
    $$
    \lim_{n \to \infty}a_n =  \lim_{n \to \infty} f(x/2^n) = f(0) = 1.
    $$
    since $f$ is continuous.

    Expanding out the recurrence, we have
    $$
    a_n = x/2^{n + 1} + x/2^{n + 2} + \ldots + x/2^{n + k} + a_{n + k}
    $$
    for all $k \geq 1.$

    Taking $k \to \infty$ on both sides, we have
    \begin{align*}
    a_n = \lim_{k \to \infty}\dfrac{x}{2^{n + 1}}\sum_{j = 0}^{k - 1} \left(
    \dfrac{1}{2} + 1
    \right)^{j} + 1 &= \dfrac{x}{2^{n + 1}} \dfrac{1}{1 - 1/2} + 1 \\
    &= \dfrac{x}{2^{n}}  +1 = f(x/2^{n})
    \end{align*}

    Letting $n = 0,$ we have $f(x) = x+ 1.$ \qed

    \item[5] Using partial fractions, we can write
    $$
    \dfrac{1}{n(n + 1)(n+2)} = \dfrac{1}{2n} - \dfrac1{n + 1} + \dfrac{1}{2(n + 2)}.
    $$
    Hence, our sum becomes
    \begin{align*}
        S &= \sum_{n = 1}^{\infty} \dfrac{1}{2n} - \dfrac{1}{n + 1} + \dfrac{1}{2(n+2)} \\
        &= \dfrac{1}2{}\sum_{n = 1}^{\infty} \dfrac{1}n{} + \dfrac{1}{2} \sum_{n = 3}^{\infty}\dfrac{1}{n} - \sum_{n = 2}^{\infty}\dfrac{1}{n} \\
        &= \dfrac{1}{2} \left(
        1 + \dfrac{1}{2}
        \right) + \sum_{n = 3}^{\infty} \dfrac{1}n{} - \dfrac{1}2{} - \sum_{n = 3}^{\infty} \dfrac{1}{n} \\
        &= \dfrac{1}{2}\left(1 + \dfrac{1}{2}\right) - \dfrac{1}{2} = \dfrac{1}4{}
    \end{align*}
\end{enumerate}

\subsection{Blue Square}
\begin{enumerate}
    \item[8] We are given a sequence $(a_n)_{n \geq 0}$ satisfying $$
    2a_n < a_{n -2} + a_{n - 1},
    $$
    meaning 
    $$
    2a_{n} + a_{n - 1} < 2a_{n -1} + a_{n - 2}.
    $$

    Letting $b_n = 2a_n + a_{n - 1},$ we get $b_n < b_{n - 1}.$ Since $b_n > 0,$ and $(b_n)_{n \geq 0}$ is monotonically decreasing, it must converge.

    Since $a_n < 2a_n + a_{n - 1} = b_n$ ($a_n > 0$), we know as a result that $a_n$ must converge too.

    \item[9] We use the substitution $x = (1-t)/(1+t),$ with differentials
    $$
    dx = -\dfrac{2}{(1+t)^2},
    $$
    yielding
    \begin{align*}
    I &= -2\int_{1}^{0} \dfrac{\ln(2) - \ln(1+t)}{(2 + 2t^2)} \, dt = \int_{0}^{1} \dfrac{\ln 2 - \ln(1+t)}{1 + t^2} \, dt \\
    &= \ln 2\int_{0}^{1} \dfrac{1}{1+t^2} \, dt - I.
    \end{align*}
    Therefore, 
    $$
    I = \dfrac{\pi\ln 2}{8}.
    $$
\end{enumerate}

\subsection{Black Diamond}
\begin{enumerate}
    \item[11] Consider the Maclaurin series for $e^x$ with infinite radius of convergence. Then, we have
    $$
        f_0(x) = \sum_{k = 0} ^{\infty} \dfrac{x^k}{k!}.
    $$
    Assume $f_n$ has Maclaurin series  $$
    \sum_{k = 0}^{\infty} F_kx^{k}.
    $$
    with infinite radius of convergence. Then, $f_{n + 1}$ will have Maclaurin series
    $$
    \sum_{k = 0}^{\infty} kF_k x^{k}
    $$
    with infinite radius of convergence, since the derivative has an infinite radius of convergence.

    Inductively, the Maclaurin series for $f_n$ satisfies 
    $$
    f_n(x) = \sum_{k = 0}^{\infty} \dfrac{k^{n} x^{k}}{k!},
    $$
    since both the base case and inductive step are satisfied. As a result, we have
    $$
    f_n(1) = \sum_{k = 0}^{\infty} \dfrac{k^n}{k!},
    $$ 
    meaning 
    \begin{align*}
    \sum_{n = 0}^{\infty} \dfrac{f_n(1)}{n!} &= \sum_{n = 0}^{\infty} \sum_{k = 0}^{\infty} \dfrac{k^n}{n! k!} \\
    &= \sum_{k = 0}^{\infty} \dfrac{1}{k!}\sum_{n = 0}^{\infty} \dfrac{k^n}{n!} \\
    &= \sum_{k = 0}^{\infty} \dfrac{e^k}{k!} = e^{e}
    \end{align*}
    from the Maclaurin series of $e^x.$
\end{enumerate}

\subsection{Reflection}
Black Diamond 11 requires dealing with Maclaurin series, and requires infinite radii of convergence for all the series involved, that requires formal statement.
\end{document}